% Compilar en Overleaf con pdfLaTeX o XeLaTeX
\documentclass[11pt,a4paper]{article}

\usepackage[utf8]{inputenc}
\usepackage[T1]{fontenc}
\usepackage[spanish,es-tabla]{babel}
\usepackage{geometry}
\usepackage{array}
\usepackage{tabularx}
\usepackage{multirow}
\usepackage{makecell}
\usepackage{ragged2e}
\usepackage{xcolor}
\usepackage{colortbl}
\usepackage{hyperref}
\usepackage{float}
\usepackage{enumitem}
\usepackage{tikz}
\usetikzlibrary{shapes,arrows.meta,positioning,fit}

\geometry{margin=2.5cm}
\hypersetup{colorlinks=true, linkcolor=blue!60!black, urlcolor=blue!60!black}

% Estilo uniforme para tablas con bordes
\definecolor{tabhead}{gray}{0.88}
\newcolumntype{L}[1]{>{\RaggedRight\arraybackslash\hspace{0pt}}p{#1}}
\newcolumntype{Y}{>{\RaggedRight\arraybackslash\hspace{0pt}}X}
\renewcommand{\arraystretch}{1.25}
\setlength{\tabcolsep}{4pt}

\title{\textbf{Investigación: IA para análisis de imágenes de electrodomésticos}\\
\large Proyecto integrador HarmoniWatts --- Cuadro comparativo}
\author{Equipo HarmoniWatts}
\date{\today}

\begin{document}
\maketitle

\begin{abstract}
Este documento compara soluciones de inteligencia artificial (local y en nube) para analizar fotografías de electrodomésticos y completar automáticamente un formulario con \emph{marca}, \emph{tipo} y \emph{potencia promedio (W)}. Se incluye Vertex AI Search como fuente de enriquecimiento, estrategias de optimización de tokens, un flujo de ampliación progresiva del catálogo tras confirmación del usuario, y una arquitectura híbrida recomendada para despliegue académico en la nube.
\end{abstract}

\section{Contexto del proyecto}

HarmoniWatts gestiona electrodomésticos con los campos: \texttt{idTipoPredefinido}, \texttt{idMarcaPredefinida}, \texttt{marcaOtro}, \texttt{nombre}, \texttt{potenciaW}, \texttt{usoSemanal} y \texttt{horarioHabitual}. El objetivo de la IA es reducir la carga manual del usuario al registrar un aparato a partir de una foto (etiqueta de potencia, frontal del equipo o manual).

\section{Flujo propuesto (Imagen $\rightarrow$ JSON $\rightarrow$ Formulario)}

\begin{figure}[h]
\centering
\begin{tikzpicture}[
  node distance=0.9cm and 1.1cm,
  box/.style={draw, rounded corners, align=center, minimum width=2.6cm, minimum height=0.85cm, font=\small},
  arrow/.style={-{Stealth[length=2mm]}, thick}
]
  \node[box, fill=blue!8] (img) {Imagen\\(Angular)};
  \node[box, fill=green!10, right=of img] (pre) {Preproceso\\(recorte, resize)};
  \node[box, fill=orange!12, right=of pre] (vlm) {Modelo\\Visión + LLM};
  \node[box, fill=yellow!15, below=of vlm] (json) {JSON estructurado\\(marca, tipo, W)};
  \node[box, fill=purple!10, left=of json] (cat) {Comparar con\\catálogo BD};
  \node[box, fill=red!8, left=of cat] (search) {Vertex AI\\Search (opc.)};
  \node[box, fill=gray!12, below=of cat] (form) {Formulario\\prellenado};
  \node[box, fill=cyan!10, below=of form] (fb) {Confirmación\\del usuario};
  \node[box, fill=lime!15, below=of fb] (enrich) {Enriquecer\\catálogo (opc.)};

  \draw[arrow] (img) -- (pre);
  \draw[arrow] (pre) -- (vlm);
  \draw[arrow] (vlm) -- (json);
  \draw[arrow] (json) -- (cat);
  \draw[arrow] (cat) -- (search);
  \draw[arrow] (search) -- (form);
  \draw[arrow] (cat) -- (form);
  \draw[arrow] (form) -- (fb);
  \draw[arrow] (fb) -- (enrich);
\end{tikzpicture}
\caption{Pipeline recomendado: visión genera JSON; se valida contra catálogo PostgreSQL; Vertex AI Search enriquece si falta información; el usuario confirma; opcionalmente se amplía el catálogo sin sobrescribir registros existentes.}
\end{figure}

\subsection{Esquema JSON objetivo}

\begin{verbatim}
{
  "marca_detectada": "Samsung",
  "modelo_detectado": "RT38K5932SL",
  "tipo_sugerido": "REFRIGERADOR",
  "potencia_w_estimada": 150,
  "confianza": 0.82,
  "fuente_potencia": "etiqueta_ocr | catalogo | busqueda_web"
}
\end{verbatim}

\section{Vertex AI Search: para qué sirve en HarmoniWatts}

Vertex AI Search es el servicio de Google Cloud que reemplaza a la API Custom Search JSON (cerrada a nuevos clientes). \textbf{No analiza imágenes}; su rol en el proyecto es complementario al modelo de visión:

\begin{itemize}[leftmargin=*]
  \item \textbf{Buscar fichas técnicas} cuando la foto no muestra claramente la potencia, el modelo exacto o la marca.
  \item \textbf{Consultar dominios acotados}: manuales de fabricantes, catálogos de retailers energéticos o documentación indexada por el equipo.
  \item \textbf{Enriquecer el JSON} con datos estructurados (marca, modelo, potencia nominal, tipo de aparato) antes de prellenar el formulario.
  \item \textbf{Reducir tokens del LLM de visión}: en lugar de pedir al modelo que ``adivine'' la potencia, se lanza una búsqueda con la marca y modelo detectados y se usa el resultado como sugerencia.
\end{itemize}

\textbf{Ejemplo de flujo:}
\begin{enumerate}[leftmargin=*]
  \item La IA de visión detecta \emph{Samsung RT38K5932SL} pero no lee la potencia en la etiqueta.
  \item El backend consulta Vertex AI Search con \texttt{Samsung RT38K5932SL potencia watts}.
  \item Se obtiene la ficha técnica y se propone \texttt{potenciaW = 150} en el formulario.
  \item El usuario confirma o corrige antes de guardar.
\end{enumerate}

En la arquitectura académica, Vertex AI Search actúa como \textbf{adaptador de enriquecimiento externo}: el backend Java envía solo texto de búsqueda (marca, modelo, tipo), nunca la imagen del usuario.

\section{Cuadro comparativo de modelos de IA}

\begin{table}[H]
\centering
\footnotesize
\setlength{\tabcolsep}{3.5pt}
\renewcommand{\arraystretch}{1.25}
\caption{Comparativa de modelos --- características y adecuación}
\label{tab:comparativa-a}
\begin{tabularx}{\textwidth}{|L{0.19\textwidth}|L{0.10\textwidth}|Y|Y|Y|}
\hline
\rowcolor{tabhead}
\textbf{Modelo} &
\textbf{Modo} &
\makecell{\textbf{Código}\\\textbf{abierto}} &
\makecell{\textbf{Visión}\\\textbf{/ OCR}} &
\makecell{\textbf{Adecuación}\\\textbf{HarmoniWatts}} \\ \hline

\textbf{GPT-4o mini} (OpenAI) &
Nube &
No &
OCR de etiquetas; identifica marca y modelo &
\textbf{Alta}: prototipo rápido con JSON mode \\ \hline

\textbf{Gemini 2.0 Flash} &
Nube &
No &
Buena relación costo/rendimiento en visión &
\textbf{Muy alta} en despliegue Google Cloud \\ \hline

\textbf{Mistral / Pixtral} &
\makecell{Nube o\\local} &
Parcial &
Lectura de etiquetas y documentos &
\textbf{Alta}: recomendación del profesor \\ \hline

\textbf{Llama 3.2 Vision 11B} &
Local &
Sí &
VQA y OCR básico &
\textbf{Alta} para demo offline \\ \hline

\textbf{Llama 3.2 Vision 90B} &
Local &
Sí &
Documentos complejos; más lento &
\textbf{Media}: excesivo para formulario básico \\ \hline

\textbf{Gemma 3} &
Local &
Sí &
Visión competente (4B--27B) &
\textbf{Media-Alta}: alternativa ligera \\ \hline

\textbf{G. Custom Search API} &
Nube &
No &
No ve imagen; enriquece JSON &
\textbf{Complementaria} (deprecated 2027) \\ \hline

\textbf{Vertex AI Search} &
Nube &
No &
Búsqueda semántica en catálogos &
\textbf{Alta}: reemplazo de Custom Search \\ \hline
\end{tabularx}
\end{table}

\begin{table}[H]
\centering
\footnotesize
\setlength{\tabcolsep}{3.5pt}
\renewcommand{\arraystretch}{1.25}
\caption{Comparativa de modelos --- costos, escalado y seguridad}
\label{tab:comparativa-b}
\begin{tabularx}{\textwidth}{|L{0.21\textwidth}|Y|Y|Y|}
\hline
\rowcolor{tabhead}
\textbf{Modelo} &
\makecell{\textbf{Costo}\\\textbf{aprox.}} &
\makecell{\textbf{Escalado}\\\textbf{horizontal}} &
\makecell{\textbf{Seguridad /}\\\textbf{privacidad}} \\ \hline

\textbf{GPT-4o mini} &
\makecell{$\sim$\$0{,}15 / 1M tokens\\entrada; $\sim$\$0{,}60 salida} &
API gestionada con autoescalado &
Datos salen del entorno; usar retención cero \\ \hline

\textbf{Gemini 2.0 Flash} &
\makecell{Capa gratuita generosa;\\$\sim$\$0{,}10--0{,}40 / 1M tokens} &
Integración nativa con GCP &
Región UE; Vertex AI con VPC \\ \hline

\textbf{Mistral / Pixtral} &
\makecell{Pixtral: $\sim$\$0{,}10 / 1M;\\Small 4: $\sim$\$0{,}15 / \$0{,}60} &
\makecell{API escalable;\\self-host con K8s (HPA)} &
\makecell{Datacenter UE;\\self-host mantiene datos locales} \\ \hline

\textbf{Llama 3.2 Vision 11B} &
\makecell{Gratis en hardware propio;\\requiere 8 GB VRAM} &
\makecell{Réplicas Ollama\\+ balanceador} &
\textbf{Máxima privacidad} (sin internet) \\ \hline

\textbf{Llama 3.2 Vision 90B} &
\makecell{$\sim$64 GB VRAM;\\GPU costosa en nube} &
Escalable pero pocas réplicas &
Privacidad total \\ \hline

\textbf{Gemma 3} &
Gratis en hardware propio &
Réplicas Ollama &
Datos no salen del servidor \\ \hline

\textbf{G. Custom Search API} &
\makecell{100 queries/día gratis;\\\$5 / 1000 queries} &
Alta; servicio stateless &
Solo envía texto de búsqueda \\ \hline

\textbf{Vertex AI Search} &
Pay-as-you-go &
Alta en GCP &
IAM, VPC y auditoría \\ \hline
\end{tabularx}
\end{table}

\section{Tabla resumida de decisión (puntuación 1--5)}

\begin{table}[H]
\centering
\footnotesize
\setlength{\tabcolsep}{3pt}
\renewcommand{\arraystretch}{1.2}
\caption{Tabla resumida de decisión (puntuación 1--5). *Custom Search no ve la imagen; puntúa enriquecimiento web. **Mistral self-host offline. El total excluye la fila con N/A.}
\begin{tabularx}{\textwidth}{|L{0.24\textwidth}|>{\centering\arraybackslash\hspace{0pt}}X|>{\centering\arraybackslash\hspace{0pt}}X|>{\centering\arraybackslash\hspace{0pt}}X|>{\centering\arraybackslash\hspace{0pt}}X|>{\centering\arraybackslash\hspace{0pt}}X|}
\hline
\rowcolor{tabhead}
\textbf{Criterio} &
\makecell{\textbf{GPT-4o}\\\textbf{mini}} &
\makecell{\textbf{Gemini}\\\textbf{Flash}} &
\makecell{\textbf{Mistral}\\\textbf{Pixtral}} &
\makecell{\textbf{Llama 3.2}\\\textbf{Vision 11B}} &
\makecell{\textbf{G. Custom}\\\textbf{Search}} \\ \hline
Precisión marca/potencia     & 4 & 4 & 4 & 3 & 2* \\ \hline
Costo académico              & 4 & 5 & 4 & 5 & 3 \\ \hline
Código abierto               & 1 & 1 & 4 & 5 & 1 \\ \hline
Funciona offline             & 1 & 1 & 3** & 5 & 1 \\ \hline
Escalado horizontal          & 5 & 5 & 4 & 3 & 5 \\ \hline
Seguridad / soberanía datos  & 2 & 3 & 4 & 5 & 3 \\ \hline
Facilidad integración Java   & 5 & 4 & 4 & 4 & 5 \\ \hline
Salida JSON estructurada     & 5 & 5 & 4 & 3 & N/A \\ \hline
\rowcolor{tabhead}
\textbf{Total (máx.\ 35)}    & \textbf{27} & \textbf{28} & \textbf{31} & \textbf{33} & \textbf{20} \\ \hline
\end{tabularx}
\end{table}

\section{Optimización de tokens (recomendaciones del profesor)}

\begin{enumerate}[leftmargin=*]
  \item \textbf{Prompt corto y schema fijo}: pedir únicamente JSON con campos del DTO \texttt{ElectrodomesticoCreateRequest}; prohibir texto libre.
  \item \textbf{Redimensionar imagen} a 512--768 px en el lado largo antes de enviarla (reduce tokens visuales hasta 60\%).
  \item \textbf{Recorte inteligente}: detectar región de etiqueta (potencia) con OpenCV antes del LLM.
  \item \textbf{Catálogo en caché}: enviar solo IDs/códigos de marcas y tipos válidos en el prompt, no todo el JSON de catálogo.
  \item \textbf{Two-step pipeline}: (1) OCR/extracción con modelo pequeño; (2) solo si confianza $< 0{,}7$, llamar modelo grande o búsqueda web.
  \item \textbf{Comparación imagen vs JSON}: el backend Java compara salida IA con \texttt{MarcaPredefinida} y \texttt{ElectrodomesticoTipoPredefinido}; fuzzy match (Levenshtein) para marcas.
  \item \textbf{Streaming desactivado} en respuestas JSON únicas (menor complejidad, misma latencia).
\end{enumerate}

\section{Arquitectura recomendada para prueba académica en nube}

\begin{table}[H]
\centering
\footnotesize
\renewcommand{\arraystretch}{1.2}
\caption{Stack alineado con HarmoniWatts (Java 21, Angular 21+, PostgreSQL).}
\begin{tabularx}{\textwidth}{|L{0.26\textwidth}|Y|}
\hline
\rowcolor{tabhead}
\textbf{Capa} & \textbf{Tecnología sugerida} \\ \hline
Frontend (Angular 21+) & Captura/subida de imagen; preview; formulario prellenado editable \\ \hline
API Gateway / Backend (Java 21) & \texttt{ElectrodomesticoVisionService}: orquesta visión, catálogo y búsqueda \\ \hline
Visión primaria (nube) & \textbf{Mistral Pixtral} o \textbf{Gemini Flash} con salida JSON \\ \hline
Visión fallback (local) & \textbf{Ollama + Llama 3.2 Vision 11B} en contenedor con GPU \\ \hline
Enriquecimiento & \textbf{Vertex AI Search}: búsqueda semántica de fichas técnicas por marca/modelo cuando la visión no alcanza \\ \hline
Persistencia & PostgreSQL: electrodomésticos del usuario, catálogo de referencia y registro de altas en catálogo \\ \hline
Escalado & Kubernetes HPA: réplicas del servicio de inferencia; cola (RabbitMQ) para picos \\ \hline
Seguridad & Keycloak (ya en proyecto), TLS, imágenes efímeras, logs sin PII, secrets en vault \\ \hline
\end{tabularx}
\end{table}

\section{Validación y confirmación del usuario}

\begin{itemize}[leftmargin=*]
  \item El formulario siempre queda editable: la IA propone valores y el usuario confirma o corrige.
  \item Si la confianza de la IA es baja, mostrar advertencia y sugerir valores del catálogo interno (\texttt{MarcaPredefinida}, \texttt{ElectrodomesticoTipoPredefinido}).
  \item Opcional: registrar correcciones del usuario en PostgreSQL para auditoría y mejora futura de prompts.
  \item Métricas de evaluación: tasa de aceptación sin edición, precisión de marca y error en potencia (W).
\end{itemize}

\section{Enriquecimiento progresivo del catálogo}

No es necesario precargar en la base de datos todos los electrodomésticos del mercado. El catálogo (\texttt{MarcaPredefinida}, \texttt{ElectrodomesticoTipoPredefinido} y, en el futuro, modelos de referencia) puede \textbf{crecer de forma orgánica} a medida que los usuarios registran aparatos con ayuda de la IA.

\subsection{Flujo tras la confirmación del usuario}

\begin{enumerate}[leftmargin=*]
  \item El usuario revisa el formulario prellenado y confirma el registro del electrodoméstico en su vivienda.
  \item El sistema evalúa si la marca, el tipo o el modelo detectado \textbf{no existen} en el catálogo de referencia.
  \item Si faltan, y el usuario lo autoriza (o mediante regla automática configurable), se crea una \textbf{entrada nueva} en el catálogo con los datos confirmados.
  \item La próxima vez que otro usuario fotografíe un aparato similar, la IA podrá resolverlo contra ese registro ya existente, reduciendo llamadas externas y tokens.
\end{enumerate}

\subsection{Reglas para evitar sobrescritura}

\begin{table}[H]
\centering
\footnotesize
\renewcommand{\arraystretch}{1.2}
\caption{Política de enriquecimiento del catálogo sin pérdida de datos.}
\begin{tabularx}{\textwidth}{|L{0.22\textwidth}|Y|}
\hline
\rowcolor{tabhead}
\textbf{Regla} & \textbf{Comportamiento} \\ \hline
Solo inserción & Si la marca o el modelo no existen, se \textbf{insertan}; nunca se modifican registros existentes de forma silenciosa. \\ \hline
Sin sobrescritura & Si ya existe una entrada con el mismo identificador (marca + modelo normalizado), \textbf{no se actualiza} la potencia ni el tipo almacenados. \\ \hline
Conflicto de datos & Si la IA propone una potencia distinta a la del catálogo, prevalece el \textbf{valor del catálogo} y se muestra al usuario la discrepancia. \\ \hline
Actualización explícita & Cualquier cambio sobre un registro existente requiere acción de \textbf{administrador} o confirmación explícita del usuario con diff visible. \\ \hline
Trazabilidad & Cada alta en catálogo guarda origen (\texttt{IA\_VISION}, \texttt{VERTEX\_SEARCH}, \texttt{USUARIO}), fecha y confianza. \\ \hline
\end{tabularx}
\end{table}

\subsection{Beneficio para el proyecto}

\begin{itemize}[leftmargin=*]
  \item Catálogo inicial mínimo (solo tipos y marcas frecuentes), escalable sin mantenimiento manual masivo.
  \item La IA aporta valor acumulativo: cada registro confirmado mejora las sugerencias futuras.
  \item Vertex AI Search cubre el ``hueco'' cuando aún no hay entrada en catálogo; la confirmación del usuario convierte esa búsqueda en conocimiento persistente.
  \item Se respeta la integridad de datos: el catálogo crece, pero no se corrompe por inferencias erróneas de la IA.
\end{itemize}

\section{Recomendación final}

\begin{enumerate}[leftmargin=*]
  \item \textbf{Prototipo académico en nube}: Mistral Pixtral o Gemini Flash + backend Java con contrato JSON estricto.
  \item \textbf{Modo offline / privacidad}: Ollama + Llama 3.2 Vision 11B en la misma VPC o máquina de laboratorio.
  \item \textbf{Enriquecimiento externo}: Vertex AI Search para fichas técnicas cuando la visión no alcance; no depender de Custom Search.
  \item \textbf{Catálogo vivo}: tras confirmación del usuario, insertar marcas/modelos nuevos sin sobrescribir registros existentes.
  \item \textbf{Siempre} dejar el formulario editable; la IA propone, el usuario confirma (human-in-the-loop).
\end{enumerate}

\section{Referencias}

\begin{itemize}[leftmargin=*]
  \item Mistral AI --- Pricing y modelos Pixtral: \url{https://mistral.ai/pricing/}
  \item Meta Llama 3.2 Vision en Ollama: \url{https://ollama.com/library/llama3.2-vision}
  \item Google Custom Search JSON API: \url{https://developers.google.com/custom-search/v1/overview}
  \item Google Vertex AI Search: \url{https://cloud.google.com/generative-ai-app-builder/docs/enterprise-search-introduction}
  \item OpenAI GPT-4o mini: \url{https://platform.openai.com/docs/models/gpt-4o-mini}
  \item Google Gemini API: \url{https://ai.google.dev/pricing}
\end{itemize}

\end{document}
